\documentclass[10pt,a4paper]{amsart}
\usepackage{amsmath,amssymb,amsthm,graphicx}
\usepackage{verbatim}
\newenvironment{reference}{\comment}{\endcomment}
%\newenvironment{reference}{}{}
\newenvironment{slogan}{\comment}{\endcomment}
\newenvironment{history}{\comment}{\endcomment}

% For commutative diagrams we use Xy-pic
\usepackage[all]{xy}

% We use 2cell for 2-commutative diagrams.
\xyoption{2cell}
% Theorem environments.
%
\theoremstyle{plain}
\newtheorem{theorem}[subsection]{Theorem}
\newtheorem{proposition}[subsection]{Proposition}
\newtheorem{lemma}[subsection]{Lemma}

\theoremstyle{definition}
\newtheorem{definition}[subsection]{Definition}
\newtheorem{example}[subsection]{Example}
\newtheorem{exercise}[subsection]{Exercise}
\newtheorem{situation}[subsection]{Situation}

\theoremstyle{remark}
\newtheorem{remark}[subsection]{Remark}
\newtheorem{remarks}[subsection]{Remarks}

\numberwithin{equation}{subsection}

% Macros
%
\def\lim{\mathop{\mathrm{lim}}\nolimits}
\def\colim{\mathop{\mathrm{colim}}\nolimits}
\def\Spec{\mathop{\mathrm{Spec}}}
\def\Hom{\mathop{\mathrm{Hom}}\nolimits}
\def\Ext{\mathop{\mathrm{Ext}}\nolimits}
\def\SheafHom{\mathop{\mathcal{H}\!\mathit{om}}\nolimits}
\def\SheafExt{\mathop{\mathcal{E}\!\mathit{xt}}\nolimits}
\def\Sch{\mathit{Sch}}
\def\Mor{\mathop{\mathrm{Mor}}\nolimits}
\def\Ob{\mathop{\mathrm{Ob}}\nolimits}
\def\Sh{\mathop{\mathit{Sh}}\nolimits}
\def\NL{\mathop{N\!L}\nolimits}
\def\CH{\mathop{\mathrm{CH}}\nolimits}
\def\proetale{{pro\text{-}\acute{e}tale}}
\def\etale{{\acute{e}tale}}
\def\QCoh{\mathit{QCoh}}
\def\Ker{\mathop{\mathrm{Ker}}}
\def\Im{\mathop{\mathrm{Im}}}
\def\Coker{\mathop{\mathrm{Coker}}}
\def\Coim{\mathop{\mathrm{Coim}}}

% Boxtimes
%
\DeclareMathSymbol{\boxtimes}{\mathbin}{AMSa}{"02}

%
% Macros for moduli stacks/spaces
%
\def\QCohstack{\mathcal{QC}\!\mathit{oh}}
\def\Cohstack{\mathcal{C}\!\mathit{oh}}
\def\Spacesstack{\mathcal{S}\!\mathit{paces}}
\def\Quotfunctor{\mathrm{Quot}}
\def\Hilbfunctor{\mathrm{Hilb}}
\def\Curvesstack{\mathcal{C}\!\mathit{urves}}
\def\Polarizedstack{\mathcal{P}\!\mathit{olarized}}
\def\Complexesstack{\mathcal{C}\!\mathit{omplexes}}
% \Pic is the operator that assigns to X its picard group, usage \Pic(X)
% \Picardstack_{X/B} denotes the Picard stack of X over B
% \Picardfunctor_{X/B} denotes the Picard functor of X over B
\def\Pic{\mathop{\mathrm{Pic}}\nolimits}
\def\Picardstack{\mathcal{P}\!\mathit{ic}}
\def\Picardfunctor{\mathrm{Pic}}
\def\Deformationcategory{\mathcal{D}\!\mathit{ef}}
\usepackage{cleveref}
\crefname{originallemma}{original lemma}{original lemmas}
\Crefname{originallemma}{Original lemma}{Original lemmas}
\makeatletter
\newlabel{algebra-lemma-Noetherian-power-ideal-kills-module}{{00L6}{source}{}{sourcepointer-1}{}}
\newlabel{algebra-lemma-Noetherian-power-ideal-kills-module@cref}{{[originallemma][0][]00L6}{[1][0][]source}{}{}{}}
\makeatother
\crefname{originallemma}{original lemma}{original lemmas}
\Crefname{originallemma}{Original lemma}{Original lemmas}
\makeatletter
\newlabel{algebra-lemma-exists-dvr}{{00PH}{source}{}{sourcepointer-2}{}}
\newlabel{algebra-lemma-exists-dvr@cref}{{[originallemma][0][]00PH}{[1][0][]source}{}{}{}}
\makeatother
\crefname{originallemma}{original lemma}{original lemmas}
\Crefname{originallemma}{Original lemma}{Original lemmas}
\makeatletter
\newlabel{algebra-lemma-ff-rings}{{00HQ}{source}{}{sourcepointer-3}{}}
\newlabel{algebra-lemma-ff-rings@cref}{{[originallemma][0][]00HQ}{[1][0][]source}{}{}{}}
\makeatother
\crefname{originallemma}{original lemma}{original lemmas}
\Crefname{originallemma}{Original lemma}{Original lemmas}
\makeatletter
\newlabel{formal-spaces-lemma-morphism-between-formal-spectra}{{0AN0}{source}{}{sourcepointer-4}{}}
\newlabel{formal-spaces-lemma-morphism-between-formal-spectra@cref}{{[originallemma][0][]0AN0}{[1][0][]source}{}{}{}}
\makeatother
\crefname{originalreference}{original reference}{original references}
\Crefname{originalreference}{Original reference}{Original references}
\makeatletter
\newlabel{lemma-faithfully-flat-rig-surjective}{{0AQX}{source}{}{sourcepointer-5}{}}
\newlabel{lemma-faithfully-flat-rig-surjective@cref}{{[originalreference][0][]0AQX}{[1][0][]source}{}{}{}}
\makeatother
\crefname{originallemma}{original lemma}{original lemmas}
\Crefname{originallemma}{Original lemma}{Original lemmas}
\makeatletter
\newlabel{more-algebra-lemma-torsion-free}{{05EA}{source}{}{sourcepointer-6}{}}
\newlabel{more-algebra-lemma-torsion-free@cref}{{[originallemma][0][]05EA}{[1][0][]source}{}{}{}}
\makeatother
\crefname{originallemma}{original lemma}{original lemmas}
\Crefname{originallemma}{Original lemma}{Original lemmas}
\makeatletter
\newlabel{properties-lemma-complement-closed-point-Jacobson}{{02IM}{source}{}{sourcepointer-7}{}}
\newlabel{properties-lemma-complement-closed-point-Jacobson@cref}{{[originallemma][0][]02IM}{[1][0][]source}{}{}{}}
\makeatother
\title{Rig-surjectivity versus faithful flatness in degree one under vanishing I-torsion}
\author{The Stacks Project Authors}
\date{stacks-project:restricted.tex:local-snapshot-20261012; GNU FDL 1.2 or later}
\begin{document}
\maketitle
\noindent\textit{Editorial note.} All delivered mathematical text is the author's own, copied from the pinned Stacks Project chapter file restricted.tex; no statement or proof step was rewritten, reordered or removed.\par\medskip
\noindent\textit{Editorial note.} The counted proof is the author's complete argument: the proof is complete in both directions\par\medskip
\noindent\textit{Editorial note.} The label lemma-flat-rig-surjective names the counted statement, and the proof environment immediately following it in the source is delivered verbatim as its proof.\par\medskip
\noindent\textit{Editorial note.} The counted statement is carried by the author's own lemma environment in the chapter source.\par\medskip
\noindent\textit{Allowed standard prerequisites (editorial).}\begin{itemize}
\item formal spectra and rig-surjectivity tested on adic maps to complete discrete valuation rings
\item the faithfully-flat case of the preceding rig-surjectivity lemma
\item vanishing of I-torsion versus I-power torsion and the Noetherian annihilator lemma
\item the Jacobson lemma on primes defining closed points of the complement of a closed point
\item embedding of a one-dimensional local domain into a discrete valuation ring and faithful flatness of maps of rings detected on spectra
\end{itemize}\medskip
\ifdefined\hypertarget\hypertarget{sourcepointer-1}{}\fi
\noindent\textit{Source reference.} Original lemma 00L6: sources/stacks-chow/tags.txt, source lines 776--776. Prerequisite pointer invoked at proof line(s) 5497; it names a standard fact of a sibling chapter that the counted argument uses, and every step of the counted proof built on it is retained verbatim in the delivered body.\par\medskip
\ifdefined\hypertarget\hypertarget{sourcepointer-2}{}\fi
\noindent\textit{Source reference.} Original lemma 00PH: sources/stacks-chow/tags.txt, source lines 931--931. Prerequisite pointer invoked at proof line(s) 5506; it names a standard fact of a sibling chapter that the counted argument uses, and every step of the counted proof built on it is retained verbatim in the delivered body.\par\medskip
\ifdefined\hypertarget\hypertarget{sourcepointer-3}{}\fi
\noindent\textit{Source reference.} Original lemma 00HQ: sources/stacks-chow/tags.txt, source lines 652--652. Prerequisite pointer invoked at proof line(s) 5521; it names a standard fact of a sibling chapter that the counted argument uses, and every step of the counted proof built on it is retained verbatim in the delivered body.\par\medskip
\ifdefined\hypertarget\hypertarget{sourcepointer-4}{}\fi
\noindent\textit{Source reference.} Original lemma 0AN0: sources/stacks-chow/tags.txt, source lines 13423--13423. Prerequisite pointer invoked at proof line(s) 5489; it names a standard fact of a sibling chapter that the counted argument uses, and every step of the counted proof built on it is retained verbatim in the delivered body.\par\medskip
\ifdefined\hypertarget\hypertarget{sourcepointer-5}{}\fi
\noindent\textit{Source reference.} Original reference 0AQX: sources/stacks-restricted/pinned-original-source.tex, source lines 5418--5418. Prerequisite pointer invoked at proof line(s) 5482; it names a standard fact of the same chapter that the counted argument uses, and every step of the counted proof built on it is retained verbatim in the delivered body.\par\medskip
\ifdefined\hypertarget\hypertarget{sourcepointer-6}{}\fi
\noindent\textit{Source reference.} Original lemma 05EA: sources/stacks-chow/tags.txt, source lines 6983--6983. Prerequisite pointer invoked at proof line(s) 5485; it names a standard fact of a sibling chapter that the counted argument uses, and every step of the counted proof built on it is retained verbatim in the delivered body.\par\medskip
\ifdefined\hypertarget\hypertarget{sourcepointer-7}{}\fi
\noindent\textit{Source reference.} Original lemma 02IM: sources/stacks-chow/tags.txt, source lines 3305--3305. Prerequisite pointer invoked at proof line(s) 5502; it names a standard fact of a sibling chapter that the counted argument uses, and every step of the counted proof built on it is retained verbatim in the delivered body.\par\medskip
\begin{lemma}
\label{lemma-flat-rig-surjective}
Let $A$ be a Noetherian ring complete with respect to an ideal $I$.
Let $B$ be an $I$-adically complete $A$-algebra. Assume that
\begin{enumerate}
\item the $I$-torsion in $A$ is $0$,
\item $A/I^n \to B/I^nB$ is flat and of finite type for all $n$.
\end{enumerate}
Then $\text{Spf}(B) \to \text{Spf}(A)$ is rig-surjective if and only
if $A/I \to B/IB$ is faithfully flat.
\end{lemma}

\begin{proof}
Faithful flatness implies rig-surjectivity by
Lemma \ref{lemma-faithfully-flat-rig-surjective}.
To prove the converse we will use without further mention that the
vanishing of $I$-torsion is equivalent to the vanishing of $I$-power torsion
(More on Algebra, Lemma \ref{more-algebra-lemma-torsion-free}).
We will also use without further mention that morphisms between
formal spectra are given by continuous maps between the corresponding
topological rings, see
Formal Spaces, Lemma \ref{formal-spaces-lemma-morphism-between-formal-spectra}.

\medskip\noindent
Assume $\text{Spf}(B) \to \text{Spf}(A)$ is rig-surjective.
Choose a maximal ideal $I \subset \mathfrak m \subset A$.
The open $U = \Spec(A_\mathfrak m) \setminus V(I_\mathfrak m)$
of $\Spec(A_\mathfrak m)$ is nonempty as the $I_\mathfrak m$-torsion of
$A_\mathfrak m$ is zero
(use Algebra, Lemma \ref{algebra-lemma-Noetherian-power-ideal-kills-module}).
Thus we can find a prime $\mathfrak q \subset A_\mathfrak m$ which defines
a point of $U$ (i.e., $IA_\mathfrak m \not \subset \mathfrak q$)
and which corresponds to a closed point
of $\Spec(A_\mathfrak m) \setminus \{\mathfrak m\}$, see
Properties, Lemma \ref{properties-lemma-complement-closed-point-Jacobson}.
Then $A_\mathfrak m/\mathfrak q$ is a dimension $1$ local domain.
Thus we can find an injective local homomorphism of local rings
$A_\mathfrak m/\mathfrak q \subset R$ where $R$ is a discrete valuation ring
(Algebra, Lemma \ref{algebra-lemma-exists-dvr}).
By construction $IR \subset \mathfrak m_R$ and we see that
$A \to R$ extends to a continuous map $A \to R^\wedge$.
Since the completion of a discrete valuation ring is a discrete
valuation ring, we see that the assumption gives a commutative
diagram of rings
$$
\xymatrix{
R' & B \ar[l] \\
R^\wedge \ar[u] & A \ar[l] \ar[u]
}
$$
Thus we find a prime ideal of $B$ lying over $\mathfrak m$. It follows
that $\Spec(B/IB) \to \Spec(A/I)$ is surjective, whence $A/I \to B/IB$
is faithfully flat
(Algebra, Lemma \ref{algebra-lemma-ff-rings}).
\end{proof}

\end{document}
